\documentclass[11pt,twoside]{article}

% ======================================================================
%  Deep Verification of the VSEPR-SIM Structural Energy Kernel (v2.7)
%  Architecture, Implementation, and Deterministic Validation
% ======================================================================

\usepackage{amsmath,amssymb,amsthm}
\usepackage{geometry}
\usepackage{booktabs}
\usepackage{longtable}
\usepackage{graphicx}
\usepackage{hyperref}
\usepackage{enumerate}
\usepackage{listings}
\usepackage{xcolor}
\usepackage{float}
\usepackage{caption}
\usepackage{subcaption}
\usepackage{fancyhdr}
\usepackage{titlesec}
\usepackage[T1]{fontenc}
\usepackage{microtype}

\geometry{letterpaper, margin=1in, headheight=14pt}

\hypersetup{
  colorlinks=true,
  linkcolor=blue!60!black,
  citecolor=green!50!black,
  urlcolor=blue!70!black
}

\lstset{
  basicstyle=\ttfamily\small,
  keywordstyle=\color{blue!80!black}\bfseries,
  commentstyle=\color{gray},
  stringstyle=\color{red!70!black},
  breaklines=true,
  columns=fullflexible,
  keepspaces=true,
  language=C++,
  frame=single,
  framesep=3pt,
  xleftmargin=6pt,
  xrightmargin=6pt,
  aboveskip=8pt,
  belowskip=8pt
}

\pagestyle{fancy}
\fancyhf{}
\fancyhead[LE,RO]{\thepage}
\fancyhead[RE]{\textit{VSEPR-SIM Deep Verification}}
\fancyhead[LO]{\textit{\nouppercase{\leftmark}}}
\renewcommand{\headrulewidth}{0.4pt}

\titleformat{\section}{\Large\bfseries}{\thesection}{1em}{}
\titleformat{\subsection}{\large\bfseries}{\thesubsection}{1em}{}

% Custom commands
\newcommand{\Angstrom}{\text{\AA}}
\newcommand{\kcalmol}{\text{kcal/mol}}
\newcommand{\Frms}{F_{\mathrm{rms}}}
\newcommand{\rmin}{r_{\mathrm{min}}}

% ======================================================================
\begin{document}

% Title page
\begin{titlepage}
\centering
\vspace*{2cm}

{\Huge\bfseries Deep Verification of the VSEPR-SIM\\
Structural Energy Kernel (v2.7)\par}

\vspace{0.8cm}
{\Large Architecture, Implementation, and\\
Deterministic Validation\par}

\vspace{2cm}

{\large Formation Engine Development Team\par}

\vspace{1cm}

{\large Version 2.7.1 --- \today\par}

\vspace{2cm}

\begin{minipage}{0.8\textwidth}
\begin{center}
\begin{tabular}{rl}
\toprule
\textbf{Verification status} & 485 / 485 PASS, 0 FAIL, 0 SKIP \\
\textbf{Deep suite}          & 256 / 256 PASS \\
\textbf{Phase executables}   & 229 / 229 PASS \\
\textbf{Canonical seed}      & \texttt{1772952709} \\
\textbf{Platform}            & GNU/Linux x86\_64 (WSL2), g++ 14.2 \\
\textbf{Repository}          & \url{https://github.com/LMSM3/VSEPR-SIM} \\
\textbf{Tag}                 & \texttt{v2.7.1} \\
\bottomrule
\end{tabular}
\end{center}
\end{minipage}

\vfill
{\small MIT License \quad|\quad C++20 \quad|\quad CMake 3.15+}
\end{titlepage}

% ======================================================================
\tableofcontents
\newpage

% ======================================================================
\section{Abstract}

VSEPR-SIM is a classical atomistic formation engine that generates molecular
and crystal structures deterministically from elemental composition and
thermodynamic boundary conditions.  This document presents the architecture,
energy model, numerical methods, and verification framework of the v2.7.1
release.

The deep verification suite comprises 256 individual checks spanning
Lennard-Jones parameter audits (42 UFF elements), Coulomb electrostatics,
dielectric screening across 19 experimental solvents, crystal geometry
validation against 30 Wyckoff reference lattices, molecular bond detection,
composite force--energy consistency, noble-gas cluster FIRE relaxation with
two-layer convergence criteria, restoring-force scans, crystal builder
geometry fidelity, randomised seeded consistency suites, NVE energy
conservation, and cross-checks against PubChem and NIST spectroscopic
reference data.

All 256 deep verification checks and 229 phased revalidation checks pass
without failure or skip.  NVE energy drift remains below
$3 \times 10^{-11}$~$\kcalmol$ on argon clusters.  FIRE minimisation
converges to $\Frms < 10^{-5}$~$\kcalmol/\Angstrom$ for all tested
noble-gas clusters (Ar$_2$--Ar$_7$).

This result establishes the first fully verified deterministic kernel
baseline for VSEPR-SIM, marking the transition from foundational
correctness recovery toward convergence tightening, expanded physical
coverage, and formation-oriented workflow integration.


% ======================================================================
\section{Project Motivation}
\label{sec:motivation}

Atomistic simulation engines fail silently when:
\begin{itemize}
  \item Force and energy terms diverge (inconsistent analytic gradient
        vs.\ numerical derivative)
  \item Geometric emitters produce incorrect crystal structures
        (wrong atom count, wrong nearest-neighbour distance)
  \item Electrostatic scaling breaks (dielectric screening not proportional
        to $\varepsilon_r$)
  \item Numerical drift accumulates (NVE total energy not conserved)
  \item Parameter tables are incomplete (fallback to default values without
        warning)
\end{itemize}

These failure modes are difficult to detect in production because individual
simulations may produce plausible-looking results while carrying systematic
errors.  VSEPR-SIM addresses this through two principles:

\begin{enumerate}
  \item \textbf{Deterministic kernel design.}  Given identical inputs, the
        engine produces bit-identical outputs.  All random number generators
        are seeded and logged.
  \item \textbf{Verification-driven development.}  Every physics module is
        exercised by a self-testing suite before integration into the
        production pipeline.  New features must pass existing checks before
        being merged.
\end{enumerate}

No fake physics---code reflects real physical modelling; visual-only
approximations are never used.


% ======================================================================
\section{Historical Development}
\label{sec:history}

\begin{table}[H]
\centering
\caption{Release timeline.}
\label{tab:versions}
\begin{tabular}{lll}
\toprule
\textbf{Tag} & \textbf{Date} & \textbf{Milestone} \\
\midrule
v2.5.0 & 2025 & Clean baseline release with 13-section methodology \\
v2.5.3 & 2025 & Heat-gated reaction control system \\
v2.6.1 & 2025 & Alpha model generative baseline \\
v2.6.2 & 2025 & Plateau investigation and root-cause analysis \\
v2.6.3 & 2025 & Empirical lookup table (118 elements, compile-time) \\
v2.7.0 & 2025 & Qt6 desktop application scaffold \\
\textbf{v2.7.1} & \textbf{2025} & \textbf{Deep verification milestone (485/485)} \\
\bottomrule
\end{tabular}
\end{table}

The progression from v2.5.0 to v2.7.1 reflects a deliberate arc: establish
the methodology (v2.5), explore generative models (v2.6), diagnose and
stabilise the kernel (v2.6.2--v2.6.3), and verify the result exhaustively
(v2.7.1).  The v2.6.2 ``plateau investigation'' was a critical turning
point---it identified four root causes (f-block degeneracy, metric mismatch,
coordinate-descent saddle points, and smoothness penalty) that motivated the
shift from generative exploration to verification-first development.


% ======================================================================
\section{System Architecture}
\label{sec:architecture}

\subsection{Pipeline Overview}

The formation engine processes structures through a layered pipeline:

\begin{center}
\begin{tabular}{cl}
\toprule
\textbf{Layer} & \textbf{Responsibility} \\
\midrule
1. Formation Priors    & Elemental identity, stoichiometry, boundary conditions \\
2. Topology Generation & Bond detection, molecular graph, crystal lattice \\
3. Energy Kernel       & LJ + Coulomb + bonded MM evaluation \\
4. Integration         & FIRE minimisation, velocity Verlet, Langevin \\
5. Structure Output    & XYZ / XYZA / XYZC with SHA-256 provenance \\
6. Verification        & Automated self-testing against empirical references \\
\bottomrule
\end{tabular}
\end{center}

\subsection{Atom Representation}

Each atom carries:
\begin{itemize}
  \item Position $\mathbf{x}_i \in \mathbb{R}^3$ (\AA)
  \item Velocity $\mathbf{v}_i \in \mathbb{R}^3$ (\AA/fs)
  \item Force $\mathbf{F}_i \in \mathbb{R}^3$ ($\kcalmol/\Angstrom$)
  \item Mass $m_i$ (amu)
  \item Atomic number $Z_i$
  \item Formal charge $q_i$ (elementary charges)
\end{itemize}

The \texttt{State} object aggregates $N$ atoms with parallel arrays for
maximum cache locality.  No atom stores pointers to other atoms---all
interactions are computed from positions at evaluation time.

\subsection{Model Interface}

All force models implement a pure-function contract:
\begin{equation}
  \texttt{eval} : \mathcal{S} \times \Theta \;\longrightarrow\; (\mathbf{F},\, \mathcal{E})
\end{equation}
where $\mathcal{S}$ is the State, $\Theta$ is the parameter set (cutoff
radius, force constants, environment context), $\mathbf{F}$ is the force
array, and $\mathcal{E}$ is the six-component energy ledger
($U_{\mathrm{LJ}}$, $U_{\mathrm{Coul}}$, $U_{\mathrm{bond}}$,
$U_{\mathrm{angle}}$, $U_{\mathrm{torsion}}$, $U_{\mathrm{external}}$).

This is a pure function: given the same State and parameters, it always
returns the same forces and energies.

\subsection{Desktop Application Layer}

The v2.7.0 release introduced a Qt6 desktop application scaffold with a
three-layer architecture:

\begin{center}
\begin{tabular}{cl}
\toprule
\textbf{Layer} & \textbf{Components} \\
\midrule
Presentation   & MainWindow, ObjectTree, PropertiesPanel, ViewportWidget \\
Bridge         & EngineAdapter (atomistic $\leftrightarrow$ Qt type conversion) \\
Engine         & Atomistic kernel (unchanged from CLI pathway) \\
\bottomrule
\end{tabular}
\end{center}

The bridge layer ensures the atomistic kernel remains independent of the GUI
framework.


% ======================================================================
\section{Energy Model}
\label{sec:energy}

\subsection{Total Energy}

The total potential energy is:
\begin{equation}
  U_{\mathrm{total}} = U_{\mathrm{LJ}} + U_{\mathrm{Coul}} + U_{\mathrm{bond}} + U_{\mathrm{angle}} + U_{\mathrm{torsion}} + U_{\mathrm{ext}}
  \label{eq:Utotal}
\end{equation}

Each term is independently computed and stored in the energy ledger.

\subsection{Lennard-Jones 12-6 Potential}

Pairwise van der Waals interactions:
\begin{equation}
  U_{\mathrm{LJ}}(r_{ij}) = 4\varepsilon_{ij} \left[
    \left(\frac{\sigma_{ij}}{r_{ij}}\right)^{12} -
    \left(\frac{\sigma_{ij}}{r_{ij}}\right)^{6}
  \right]
  \label{eq:LJ}
\end{equation}

Parameters are drawn from the Universal Force Field (UFF) of
Rapp\'{e} et al.\ (1992).  VSEPR-SIM v2.7.1 includes 42 elements with
verified $\sigma$ and $\varepsilon$ values.

Lorentz--Berthelot mixing rules:
\begin{align}
  \sigma_{ij} &= \frac{\sigma_i + \sigma_j}{2} \\
  \varepsilon_{ij} &= \sqrt{\varepsilon_i \cdot \varepsilon_j}
\end{align}

The equilibrium distance is:
\begin{equation}
  \rmin = 2^{1/6} \sigma
\end{equation}

\subsection{Coulomb Electrostatics}

Point-charge electrostatics with distance-dependent dielectric screening:
\begin{equation}
  U_{\mathrm{Coul}}(r_{ij}) = \frac{k_e \, q_i \, q_j}{\varepsilon_r \, r_{ij}}
  \label{eq:coulomb}
\end{equation}

where $k_e = 332.0637$~$\kcalmol \cdot \Angstrom / e^2$ and $\varepsilon_r$
is the relative dielectric constant.

The dielectric model supports:
\begin{itemize}
  \item Near-vacuum ($\varepsilon_r = 1.0$) for gas-phase systems
  \item Explicit solvent dielectrics from CRC Handbook, 104th Edition
        (19 solvents from hexane $\varepsilon_r = 1.88$ to water
        $\varepsilon_r = 78.4$)
\end{itemize}

\subsection{Bonded Terms}

\textbf{Bond stretch} (harmonic):
\begin{equation}
  U_{\mathrm{bond}} = \frac{1}{2} k_r (r - r_0)^2
\end{equation}

\textbf{Angle bend} (harmonic):
\begin{equation}
  U_{\mathrm{angle}} = \frac{1}{2} k_\theta (\theta - \theta_0)^2
\end{equation}

\textbf{Torsion} (cosine series):
\begin{equation}
  U_{\mathrm{torsion}} = \sum_n \frac{V_n}{2} \bigl[1 + \cos(n\phi - \gamma_n)\bigr]
\end{equation}

Force constants are derived from UFF where available.

\subsection{Compile-Time Empirical Database}

The \texttt{empirical\_reference.hpp} header provides compile-time arrays:

\begin{center}
\begin{tabular}{lr}
\toprule
\textbf{Table} & \textbf{Entries} \\
\midrule
Bond lengths (experimental) & 57 \\
Bond angles & 25 \\
Crystal lattice parameters & 30 \\
LJ parameters (UFF) & 42 \\
Diatomic spectroscopic constants (NIST) & 10 \\
Solvent dielectrics (CRC 104) & 19 \\
Shannon ionic radii & 15 \\
Noble gas dimer parameters & 5 \\
\bottomrule
\end{tabular}
\end{center}

All values carry tolerance bounds and literature source annotations.


% ======================================================================
\section{Structural Systems}
\label{sec:structural}

\subsection{Bond Detection}

Bonds are detected by covalent-radius overlap:
\begin{equation}
  r_{ij} < f \cdot (R_{\mathrm{cov},i} + R_{\mathrm{cov},j})
\end{equation}
where $f = 1.15$ is the tolerance factor.  The algorithm runs in
$O(N^2)$ for small molecules and uses cell lists for larger systems.

Suite~E validates this against 9 canonical molecules (H$_2$ through CH$_4$),
confirming exact bond counts.

\subsection{Crystal Emitter}

Crystal structures are built from unit cell definitions with fractional
coordinates, space group metadata, and lattice parameters:

\begin{center}
\begin{tabular}{llrl}
\toprule
\textbf{Preset} & \textbf{Type} & \textbf{$a$ (\AA)} & \textbf{Space group} \\
\midrule
Cu  & FCC      & 3.615 & Fm$\overline{3}$m (225) \\
Fe  & BCC      & 2.870 & Im$\overline{3}$m (229) \\
NaCl & Rocksalt & 5.640 & Fm$\overline{3}$m (225) \\
Si  & Diamond  & 5.431 & Fd$\overline{3}$m (227) \\
Al  & FCC      & 4.050 & Fm$\overline{3}$m (225) \\
MgO & Rocksalt & 4.212 & Fm$\overline{3}$m (225) \\
CsCl & CsCl    & 4.123 & Pm$\overline{3}$m (221) \\
\bottomrule
\end{tabular}
\end{center}

Supercell replication is exact: an $n_a \times n_b \times n_c$ supercell
contains $n_a \cdot n_b \cdot n_c \cdot Z_{\mathrm{cell}}$ atoms, verified
by Suite~L for random $(n_a, n_b, n_c)$ triplets.

\subsection{Crystal Builder Geometry Fidelity}

Suite~Q verifies the builder output for four preset crystals:
\begin{itemize}
  \item \textbf{Atom count}: $N = 8 \times Z_{\mathrm{cell}}$ for all
        $2 \times 2 \times 2$ supercells
  \item \textbf{Nearest-neighbour distance}: within 3\% of the Wyckoff
        reference value
\end{itemize}

Note: FIRE relaxation is \emph{not} applied to crystal supercells in the
current suite because UFF LJ $\rmin$ exceeds the actual bond length for
metals (Cu, Fe) and covalent crystals (Si), placing atoms inside the LJ
repulsive wall.  For NaCl, the Coulomb--LJ equilibrium differs from the
physical crystal geometry.  This is documented as a known scope boundary.


% ======================================================================
\section{Numerical Methods}
\label{sec:numerical}

\subsection{FIRE Minimisation}

The Fast Inertial Relaxation Engine (Bitzek et al., 2006) is the primary
structure optimiser.  The implementation uses Euler integration with a
velocity kick after power-negative resets:

\begin{enumerate}
  \item Evaluate forces $\mathbf{F}_i$ at current positions
  \item Compute power $P = \sum_i \mathbf{v}_i \cdot \mathbf{F}_i$
  \item Mix velocities: $\mathbf{v}_i \leftarrow (1-\alpha)\mathbf{v}_i +
        \alpha |\mathbf{v}| \hat{\mathbf{F}}_i$
  \item If $P > 0$: increment positive-power counter; increase
        $\Delta t$ and decrease $\alpha$
  \item If $P \leq 0$: kick $\mathbf{v}_i = (\Delta t / |\mathbf{F}|)\,
        \mathbf{F}_i$; decrease $\Delta t$; reset $\alpha$
  \item Update positions: $\mathbf{x}_i \leftarrow \mathbf{x}_i +
        \mathbf{v}_i \Delta t$
\end{enumerate}

\subsubsection{Euler--FIRE Deadlock Correction}

A deadlock condition was identified and corrected during v2.7.1 development.
In the original implementation, the $P < 0$ branch set
$\mathbf{v}_i = \mathbf{0}$.  Because the integrator is pure Euler
($\mathbf{x} \leftarrow \mathbf{x} + \mathbf{v}\Delta t$, no acceleration
term), zero velocity means zero displacement, which yields
$\Delta U = 0$ on the next evaluation, falsely triggering the
$\varepsilon_U$ convergence criterion.

The correction replaces the zero-velocity reset with a kick along the
normalised force direction:
\begin{equation}
  \mathbf{v}_i \leftarrow \frac{\Delta t}{|\mathbf{F}|} \, \mathbf{F}_i
  \qquad \text{(after $P < 0$ event)}
\end{equation}
and conditions the $\varepsilon_U$ check on $\varepsilon_U > 0$.

\subsection{Velocity Verlet (NVE)}

The standard four-step symplectic integrator:
\begin{align}
  \mathbf{v}_i(t + \tfrac{\Delta t}{2}) &= \mathbf{v}_i(t) +
    \frac{\mathbf{F}_i(t)}{m_i} \cdot \text{ACC\_CONV} \cdot \frac{\Delta t}{2} \\
  \mathbf{x}_i(t + \Delta t) &= \mathbf{x}_i(t) +
    \mathbf{v}_i(t + \tfrac{\Delta t}{2}) \cdot \Delta t \\
  \mathbf{F}_i(t + \Delta t) &= -\nabla_{\mathbf{x}_i} U\bigl(\mathbf{X}(t+\Delta t)\bigr) \\
  \mathbf{v}_i(t + \Delta t) &= \mathbf{v}_i(t + \tfrac{\Delta t}{2}) +
    \frac{\mathbf{F}_i(t+\Delta t)}{m_i} \cdot \text{ACC\_CONV} \cdot \frac{\Delta t}{2}
\end{align}

Suite~M verifies energy conservation to $< 3 \times 10^{-11}$~$\kcalmol$ on
Ar$_3$--Ar$_5$ clusters.

\subsection{Langevin Thermostat (NVT)}

BAOAB splitting for the Langevin equation, providing canonical (NVT)
sampling with a tunable friction coefficient $\gamma$.


% ======================================================================
\section{Verification Framework}
\label{sec:framework}

\subsection{Design Philosophy}

The verification framework is built on five principles:
\begin{enumerate}
  \item \textbf{Self-contained executables.}  Each phase is a standalone C++
        program that prints PASS/FAIL.  No external test frameworks.
  \item \textbf{Empirical grounding.}  Reference values come from published
        literature (Rapp\'{e}, Wyckoff, CRC, NIST) with explicit tolerances.
  \item \textbf{Reproducibility.}  Random suites are seeded; the seed is
        printed in every report.  Re-running with \texttt{--seed N}
        reproduces the exact result.
  \item \textbf{Layered severity.}  Suite~G demonstrates two-layer checks:
        Layer~1 (descent) is always required; Layer~2 (convergence quality)
        adds $\Frms$, overlap, and stagnation guards.
  \item \textbf{Tiered workflow.}  Quick local runs skip network-dependent
        Suite~P; full release verification includes it.
\end{enumerate}

\subsection{Phased Revalidation (Phases 1--14)}

\begin{center}
\begin{tabular}{lrl}
\toprule
\textbf{Phase} & \textbf{Checks} & \textbf{Status} \\
\midrule
1 (kernel)      & 11  & PASS \\
2 (structural)  & 15  & PASS \\
3 (relaxation)  & 17  & PASS \\
4 (presets)     & 52  & PASS \\
5 (crystals)    & 13  & PASS \\
6/7 (sessions)  & 43  & PASS \\
8--10 (env)     & 28  & PASS \\
11--13 (comp)   & 26  & PASS \\
14 (milestones) & 24  & PASS \\
\midrule
\textbf{Total}  & \textbf{229} & \textbf{ALL PASS} \\
\bottomrule
\end{tabular}
\end{center}


% ======================================================================
\section{Deep Verification Suite}
\label{sec:deep}

The deep verification suite (\texttt{apps/deep\_verification.cpp}, 1213
lines) exercises the kernel across 16 suites.

\subsection{Suite Catalogue}

\begin{longtable}{clrl}
\toprule
\textbf{Suite} & \textbf{Description} & \textbf{Checks} & \textbf{Type} \\
\midrule
\endfirsthead
\toprule
\textbf{Suite} & \textbf{Description} & \textbf{Checks} & \textbf{Type} \\
\midrule
\endhead
A & LJ parameter audit (42 UFF elements vs Rapp\'{e}) & 42 & Static \\
B & Homonuclear dimer sweeps (energy, $\rmin$, smoothness) & 20 & Static \\
C & Ion-pair Coulomb + force consistency & 15 & Static \\
D & Crystal geometry vs Wyckoff (30 structures) & 90 & Static \\
E & Molecule bond detection (9 molecules) & 9 & Static \\
F & Composite force--energy (H$_2$O, NH$_3$, CH$_4$, Ar$_2$) & 4 & Static \\
G & Noble gas cluster FIRE (Ar$_2$--Ar$_7$, two-layer) & 30 & Static \\
H & Dielectric screening (19 CRC 104 solvents) & 19 & Static \\
N & Restoring-force scan (He$_2$--Xe$_2$) & 10 & Static \\
Q & Crystal builder fidelity ($2\times2\times2$ supercell) & 8 & Static \\
\midrule
I & Random LJ force--energy ($N_{\mathrm{rand}}$ pairs) & 1 & Randomised \\
J & Random Coulomb $1/r$ law ($N_{\mathrm{rand}}$ pairs) & 2 & Randomised \\
K & Random dielectric ratio ($N_{\mathrm{rand}}$ samples) & 1 & Randomised \\
L & Random supercell atom count (6 crystals) & 6 & Randomised \\
M & NVE energy conservation (Ar$_3$--Ar$_5$) & 6 & Randomised \\
\midrule
P & Downloaded empirical refs (PubChem + NIST) & 53 & External \\
\midrule
  & \textbf{Total} & \textbf{256} & \\
\bottomrule
\end{longtable}

\subsection{Suite A: LJ Parameter Audit}

For each of 42 UFF elements, verify:
\begin{itemize}
  \item $\sigma > 0$
  \item $\varepsilon > 0$
  \item $\rmin = 2^{1/6} \sigma$ (analytic identity)
\end{itemize}

\subsection{Suite B: Homonuclear Dimer Sweeps}

For noble gases and selected elements, sweep $r \in [0.8\sigma, 3\sigma]$
and verify:
\begin{itemize}
  \item Energy minimum occurs at $r = \rmin$ (within tolerance)
  \item Well depth equals $-\varepsilon$
  \item Energy curve is smooth (no discontinuities from switching)
\end{itemize}

\subsection{Suite D: Crystal Geometry}

For 30 reference crystals, build the unit cell and verify:
\begin{itemize}
  \item Nearest-neighbour distance matches Wyckoff within tolerance
  \item Coordination number is correct
  \item Atoms per cell equals $Z_{\mathrm{cell}}$
\end{itemize}

\subsection{Suite G: Noble Gas Cluster FIRE (Two-Layer)}

Starting geometry: Fibonacci sphere with radius
$R = \rmin \sqrt{N}/3.0$, placing all pairs at $> \rmin$ (attractive well).

\textbf{Layer 1} (descent):
\begin{itemize}
  \item Energy is finite
  \item $U_{\mathrm{final}} < U_0$ (energy decreased)
\end{itemize}

\textbf{Layer 2} (convergence quality):
\begin{itemize}
  \item $\Frms < 0.10$~$\kcalmol/\Angstrom$
  \item $r_{\mathrm{nn}} > 0.55 \cdot \rmin$ (no atom overlap)
  \item $\Delta U / N < 10^{-4}$ (not stagnating)
\end{itemize}

\subsection{Suite N: Restoring-Force Scan}

For each noble gas (He, Ne, Ar, Kr, Xe), place a homonuclear dimer at
$r = \rmin + \delta$ for $\delta \in \{-1.5, -1.0, -0.6, -0.3, +0.3,
+0.6, +1.0, +1.5\}$~\AA\ and verify:
\begin{itemize}
  \item Force direction is restoring: $F_x \cdot \delta > 0$
  \item $\rmin$ is the energy minimum: $U(\rmin + \delta) > U(\rmin)$ for
        all $\delta \neq 0$
\end{itemize}

\subsection{Suite P: Downloaded Empirical References}

The \texttt{scripts/fetch\_empirical.py} downloader retrieves:
\begin{itemize}
  \item 25 bond lengths from PubChem 3D conformers
  \item 20 diatomic spectroscopic $r_e$ from NIST WebBook
  \item 15 solvent dielectrics from CRC Handbook, 104th Edition
\end{itemize}

Suite~P cross-checks these against the hardcoded database and runs
solvent-specific Coulomb screening tests.  The downloader includes a
\texttt{DiskSpaceGuard} with backup/rollback protection and optional-package
detection.


% ======================================================================
\section{Verification Results}
\label{sec:results}

\subsection{Summary}

\begin{center}
\fbox{\parbox{0.7\textwidth}{\centering\large
\textbf{256 / 256 PASS \quad 0 FAIL \quad 0 SKIP}\\[4pt]
Seed: \texttt{1772952709} \quad Platform: g++ 14.2, WSL2 x86\_64}}
\end{center}

\subsection{Suite G --- FIRE Convergence Metrics}

\begin{table}[H]
\centering
\caption{FIRE relaxation of noble gas clusters from Fibonacci sphere
starting geometry.  All clusters converge to $\Frms < 10^{-5}$~$\kcalmol/\Angstrom$.}
\label{tab:fire}
\begin{tabular}{lrrrrr}
\toprule
\textbf{Cluster} & $U_0$ & $U_f$ & $\Frms$ & \textbf{Steps} & $r_{\mathrm{nn}}$ (\AA) \\
\midrule
Ar$_2$ & $-0.026$ & $-0.238$ & $9.98 \times 10^{-6}$ & 603 & 3.817 \\
Ar$_3$ & $-0.696$ & $-0.714$ & $5.35 \times 10^{-6}$ & 308 & 3.817 \\
Ar$_4$ & $-1.219$ & $-1.428$ & $4.78 \times 10^{-6}$ & 919 & 3.816 \\
Ar$_5$ & $-1.675$ & $-2.167$ & $8.72 \times 10^{-6}$ & 1269 & 3.808 \\
Ar$_6$ & $-2.077$ & $-3.025$ & $4.97 \times 10^{-6}$ & 2075 & 3.799 \\
Ar$_7$ & $-2.462$ & $-3.928$ & $8.55 \times 10^{-6}$ & 2491 & 3.792 \\
\bottomrule
\end{tabular}
\end{table}

\subsection{Suite M --- NVE Energy Conservation}

\begin{table}[H]
\centering
\caption{NVE velocity-Verlet energy conservation on argon clusters.
Maximum drift over 1000 steps at $\Delta t = 1$~fs.}
\label{tab:nve}
\begin{tabular}{lrr}
\toprule
\textbf{System} & $E_0$ ($\kcalmol$) & \textbf{Max drift} \\
\midrule
Ar$_3$ & 1.167 & $2.91 \times 10^{-12}$ \\
Ar$_4$ & 1.817 & $5.27 \times 10^{-12}$ \\
Ar$_5$ & 0.314 & $2.26 \times 10^{-11}$ \\
\bottomrule
\end{tabular}
\end{table}

\subsection{Suite I/J/K --- Randomised Consistency}

\begin{table}[H]
\centering
\caption{Worst-case relative error across 500 randomised samples per suite.}
\label{tab:random}
\begin{tabular}{llr}
\toprule
\textbf{Suite} & \textbf{Property tested} & \textbf{Worst error} \\
\midrule
I & LJ force--energy consistency & $< 10^{-3}$ \\
J & Coulomb $U = k q_1 q_2 / r$ & $< 10^{-6}$ \\
K & Dielectric $U_{\mathrm{vac}} / U_\varepsilon = \varepsilon_r$ & $< 10^{-6}$ \\
\bottomrule
\end{tabular}
\end{table}


% ======================================================================
\section{Demonstration Cases}
\label{sec:demos}

\subsection{Tier 1: Verification Anchor Systems}

These systems are the primary validation targets, present in both the static
empirical database and the PubChem/NIST downloaded cross-check.

\begin{table}[H]
\centering
\caption{Tier 1 verification anchor molecules.}
\label{tab:tier1}
\begin{tabular}{llrrl}
\toprule
\textbf{Molecule} & \textbf{Bond} & \textbf{$r_{\mathrm{ref}}$ (\AA)} & \textbf{Bonds} & \textbf{Source} \\
\midrule
H$_2$  & H--H & 0.741 & 1 & NIST \\
H$_2$O & O--H & 0.958 & 2 & PubChem \\
NH$_3$ & N--H & 1.012 & 3 & PubChem \\
CH$_4$ & C--H & 1.087 & 4 & PubChem \\
HF     & H--F & 0.917 & 1 & NIST \\
HCl    & H--Cl& 1.275 & 1 & NIST \\
CO$_2$ & C=O  & 1.160 & 2 & PubChem \\
NaCl   & Na--Cl & 2.361 & -- & NIST \\
\bottomrule
\end{tabular}
\end{table}

\subsection{Tier 2: Intermediate Molecular Demonstrations}

These molecules exercise bond detection across functional groups, ring
structures, and competing bond environments.

\begin{table}[H]
\centering
\caption{Tier 2 molecular demonstrations (PubChem 3D cross-checked).}
\label{tab:tier2}
\begin{tabular}{llrl}
\toprule
\textbf{Molecule} & \textbf{Bond tested} & \textbf{$r$ (\AA)} & \textbf{Feature} \\
\midrule
Ethane (C$_2$H$_6$) & C--C & 1.536 & Single bond \\
Ethylene (C$_2$H$_4$) & C=C & 1.339 & Double bond \\
Acetylene (C$_2$H$_2$) & C$\equiv$C & 1.203 & Triple bond \\
Benzene (C$_6$H$_6$) & C--C & 1.397 & Aromatic ring \\
Methanol (CH$_3$OH) & C--O & 1.428 & Alcohol \\
Formaldehyde (CH$_2$O) & C=O & 1.206 & Carbonyl \\
Formic acid (HCOOH) & C--O & 1.202 & Carboxyl \\
Acetone & C=O & 1.213 & Ketone \\
HCN & C$\equiv$N & 1.156 & Nitrile \\
SO$_2$ & S=O & 1.432 & Sulphur oxide \\
CS$_2$ & C=S & 1.554 & Thiocarbonyl \\
\bottomrule
\end{tabular}
\end{table}

\subsection{Tier 3: Crystal Demonstrations}

\begin{table}[H]
\centering
\caption{Crystal preset verification.  All atom counts and nearest-neighbour
distances match Wyckoff references within 3\%.}
\label{tab:crystals}
\begin{tabular}{llrlrl}
\toprule
\textbf{Crystal} & \textbf{Type} & \textbf{$a$ (\AA)} & \textbf{$r_{\mathrm{nn}}$ (\AA)} & \textbf{CN} & \textbf{$Z$} \\
\midrule
Cu  & FCC      & 3.615 & 2.556 & 12 & 4 \\
Al  & FCC      & 4.050 & 2.863 & 12 & 4 \\
Fe  & BCC      & 2.870 & 2.485 & 8  & 2 \\
NaCl & Rocksalt & 5.640 & 2.820 & 6  & 8 \\
Si  & Diamond  & 5.431 & 2.352 & 4  & 8 \\
MgO & Rocksalt & 4.212 & 2.106 & 6  & 8 \\
CsCl & CsCl    & 4.123 & 3.571 & 8  & 2 \\
\bottomrule
\end{tabular}
\end{table}

\subsection{Tier 4: Dielectric Screening Demonstrations}

\begin{table}[H]
\centering
\caption{Solvent dielectric screening verification.  All ratios
$U_{\mathrm{vac}} / U_\varepsilon = \varepsilon_r$ to $< 10^{-6}$ relative error.}
\label{tab:solvents}
\begin{tabular}{lrl}
\toprule
\textbf{Solvent} & $\varepsilon_r$ & \textbf{Source} \\
\midrule
Hexane          & 1.88  & CRC 104 \\
Benzene         & 2.27  & CRC 104 \\
Toluene         & 2.38  & CRC 104 \\
Diethyl ether   & 4.27  & CRC 104 \\
Chloroform      & 4.81  & CRC 104 \\
THF             & 7.58  & CRC 104 \\
Dichloromethane & 8.93  & CRC 104 \\
Acetone         & 20.70 & CRC 104 \\
Ethanol         & 24.85 & CRC 104 \\
Methanol        & 32.66 & CRC 104 \\
Acetonitrile    & 35.69 & CRC 104 \\
DMF             & 36.70 & CRC 104 \\
Glycerol        & 46.53 & CRC 104 \\
DMSO            & 46.70 & CRC 104 \\
Water           & 78.40 & CRC 104 \\
\bottomrule
\end{tabular}
\end{table}


% ======================================================================
\section{Limitations}
\label{sec:limitations}

This section states the known boundaries of the current kernel.

\begin{enumerate}
  \item \textbf{Metal/covalent crystal relaxation.}  UFF LJ $\rmin$ exceeds
        the actual bond length for Cu ($\rmin = 3.92$~\AA\ vs
        $r_{\mathrm{nn}} = 2.56$~\AA), Fe, and Si.  FIRE relaxation drives
        these crystals away from physical equilibrium.
  \item \textbf{No dipole interactions.}  The current model treats atoms as
        point charges; induced dipoles and polarisation are not implemented.
  \item \textbf{Incomplete charge equilibration.}  QEq / Mulliken charge
        equilibration is present in the codebase but not yet exercised by
        the verification suite.
  \item \textbf{No reactive chemistry.}  Bond formation and breaking during
        dynamics is not supported; topology is fixed at initialisation.
  \item \textbf{PBC not exercised in Suites G/N/Q.}  Noble gas clusters and
        crystal builder tests run without periodic boundary conditions.
  \item \textbf{Bonded force field coverage.}  Harmonic bond/angle terms are
        present but not yet systematically verified against experimental
        equilibrium geometries.
  \item \textbf{Suite G tests local descent only.}  The Fibonacci-sphere
        starting geometry finds a local minimum, not necessarily the global
        LJ minimum for $N \geq 4$.
\end{enumerate}


% ======================================================================
\section{Future Work}
\label{sec:future}

\begin{enumerate}
  \item \textbf{Bond/angle perturbation scans on polyatomics.}  Extend
        Suite~N to H$_2$O, NH$_3$ with bonded force fields.
  \item \textbf{PBC crystal stability.}  NVE dynamics in periodic supercells.
  \item \textbf{Defect perturbation.}  Vacancy and interstitial tests in NaCl
        and Cu supercells; verify topology preservation.
  \item \textbf{QEq charge equilibration.}  Verify convergence and Coulomb
        consistency under charge redistribution.
  \item \textbf{Basin multiplicity.}  Multi-start FIRE on Ar$_7$ to map the
        LJ local-minima landscape.
  \item \textbf{Dipole and polarisation.}  Point-dipole interactions and
        self-consistent field polarisation.
  \item \textbf{Reactive MD.}  Bond formation/breaking with energy-gated
        topology updates.
  \item \textbf{Polymorph discovery.}  Automated crystal structure search
        using formation priors and basin mapping.
\end{enumerate}


% ======================================================================
\section{Conclusion}
\label{sec:conclusion}

The deterministic kernel of VSEPR-SIM has been verified for:

\begin{itemize}
  \item Lennard-Jones interactions (42 UFF elements, force--energy
        consistency $< 10^{-3}$ relative error)
  \item Coulomb electrostatics (exact $1/r$ scaling, dielectric screening
        across 19 solvents)
  \item Crystal geometry (30 Wyckoff reference lattices, all within
        tolerance)
  \item Molecular bond detection (9 molecules, exact counts)
  \item FIRE relaxation (Ar$_2$--Ar$_7$, $\Frms < 10^{-5}$, no deadlock)
  \item NVE energy conservation ($< 3 \times 10^{-11}$~$\kcalmol$ drift)
  \item Empirical cross-check (25 PubChem bonds, 20 NIST diatomics)
\end{itemize}

All 485 verification checks pass without failure.  The result is
reproducible via seeded random suites and archived verbose output.

This establishes the first fully verified deterministic kernel baseline for
VSEPR-SIM, completing the transition from foundational correctness recovery
to a stable platform for convergence tightening, expanded physical coverage,
and formation-oriented workflow integration.

\vspace{12pt}
\noindent\textit{Reproduce this result:}
\begin{lstlisting}[language=bash,frame=none]
wsl bash scripts/run_all_verification.sh --full
wsl bash -c "cd wsl-build && ./deep_verification --seed 1772952709 --rand-iters 500"
\end{lstlisting}


% ======================================================================
\appendix

\section{UFF Lennard-Jones Parameters (42 Elements)}
\label{app:uff}

\begin{longtable}{lrrr}
\toprule
\textbf{Element} & $\sigma$ (\AA) & $\varepsilon$ ($\kcalmol$) & $\rmin$ (\AA) \\
\midrule
\endfirsthead
\toprule
\textbf{Element} & $\sigma$ (\AA) & $\varepsilon$ ($\kcalmol$) & $\rmin$ (\AA) \\
\midrule
\endhead
H  & 2.886 & 0.044 & 3.239 \\
He & 2.362 & 0.056 & 2.651 \\
Li & 2.451 & 0.025 & 2.751 \\
Be & 2.745 & 0.085 & 3.081 \\
B  & 3.637 & 0.180 & 4.082 \\
C  & 3.431 & 0.105 & 3.851 \\
N  & 3.261 & 0.069 & 3.660 \\
O  & 3.118 & 0.060 & 3.500 \\
F  & 2.997 & 0.050 & 3.364 \\
Ne & 3.243 & 0.042 & 3.640 \\
Na & 3.328 & 0.030 & 3.735 \\
Mg & 3.021 & 0.111 & 3.391 \\
Al & 4.499 & 0.505 & 5.049 \\
Si & 3.826 & 0.402 & 4.295 \\
P  & 3.695 & 0.305 & 4.147 \\
S  & 3.595 & 0.274 & 4.035 \\
Cl & 3.516 & 0.227 & 3.947 \\
Ar & 3.400 & 0.238 & 3.817 \\
K  & 4.241 & 0.035 & 4.760 \\
Ca & 3.399 & 0.238 & 3.816 \\
Ti & 2.829 & 0.017 & 3.175 \\
V  & 2.801 & 0.016 & 3.143 \\
Cr & 2.693 & 0.015 & 3.022 \\
Mn & 2.638 & 0.013 & 2.960 \\
Fe & 2.594 & 0.013 & 2.912 \\
Co & 2.559 & 0.014 & 2.872 \\
Ni & 2.525 & 0.015 & 2.834 \\
Cu & 3.114 & 0.005 & 3.495 \\
Zn & 2.462 & 0.124 & 2.763 \\
Ga & 3.905 & 0.415 & 4.384 \\
Ge & 3.813 & 0.379 & 4.280 \\
As & 3.769 & 0.309 & 4.230 \\
Se & 3.746 & 0.291 & 4.204 \\
Br & 3.732 & 0.251 & 4.189 \\
Kr & 3.689 & 0.220 & 4.141 \\
Rb & 4.114 & 0.040 & 4.617 \\
Sr & 3.641 & 0.235 & 4.088 \\
Mo & 2.719 & 0.056 & 3.052 \\
Pd & 2.583 & 0.048 & 2.899 \\
Ag & 2.805 & 0.036 & 3.149 \\
I  & 4.009 & 0.339 & 4.500 \\
Xe & 3.924 & 0.332 & 4.405 \\
\bottomrule
\end{longtable}

\noindent Source: Rapp\'{e} et al., \textit{J.\ Am.\ Chem.\ Soc.}
\textbf{114}, 10024 (1992).  All values verified by Suite~A.


\section{Empirical Bond Length Database}
\label{app:bonds}

57 reference bond lengths from NIST CCCBDB, Wyckoff crystal data, and
PubChem 3D conformers.  See \texttt{atomistic/core/empirical\_reference.hpp}
for the complete table with source annotations and tolerance bounds.


\section{CLI and Verification Usage}
\label{app:cli}

\begin{lstlisting}[language=bash,frame=none]
# Quick verification (no network)
wsl bash scripts/run_all_verification.sh

# Full verification with PubChem/NIST download
wsl bash scripts/run_all_verification.sh --full

# Deep verification with specific seed
./wsl-build/deep_verification --seed 1772952709 --verbose

# Fetch empirical references
python3 scripts/fetch_empirical.py --offline  # fallback only
python3 scripts/fetch_empirical.py            # live PubChem + NIST
python3 scripts/fetch_empirical.py --verbose  # show each fetch
\end{lstlisting}


% ======================================================================
\section*{References}

\begin{enumerate}
  \item A.~K.~Rapp\'{e}, C.~J.~Casewit, K.~S.~Colwell, W.~A.~Goddard~III,
        and W.~M.~Skiff, ``UFF, a Full Periodic Table Force Field for
        Molecular Mechanics and Molecular Dynamics Simulations,''
        \textit{J.\ Am.\ Chem.\ Soc.} \textbf{114}, 10024--10035 (1992).

  \item E.~Bitzek, P.~Koskinen, F.~G\"{a}hler, M.~Moseler, and P.~Gumbsch,
        ``Structural Relaxation Made Simple,''
        \textit{Phys.\ Rev.\ Lett.} \textbf{97}, 170201 (2006).

  \item W.~Kabsch, ``A Solution for the Best Rotation to Relate Two Sets of
        Vectors,'' \textit{Acta Cryst.} \textbf{A32}, 922--923 (1976).

  \item R.~W.~G.~Wyckoff, \textit{Crystal Structures}, 2nd~ed.
        (Interscience, New York, 1963).

  \item D.~R.~Lide, ed., \textit{CRC Handbook of Chemistry and Physics},
        104th~ed.\ (CRC Press, Boca Raton, 2023).

  \item National Institute of Standards and Technology, \textit{NIST
        Chemistry WebBook}, \url{https://webbook.nist.gov/}.

  \item National Center for Biotechnology Information, \textit{PubChem},
        \url{https://pubchem.ncbi.nlm.nih.gov/}.

  \item R.~D.~Shannon, ``Revised Effective Ionic Radii and Systematic
        Studies of Interatomic Distances in Halides and Chalcogenides,''
        \textit{Acta Cryst.} \textbf{A32}, 751--767 (1976).
\end{enumerate}

\end{document}
